% !TeX root = ../../../poisson-cubic.tex
\section{非楕円型ブロックを持つ固定Poisson成分}\label{ne:examples}
3変数制限が特異でも、中心変数との混合項により超曲面全体は滑らかになり得る。
ここでは、節点型または尖点型の3変数制限を一般点に持つ固定Poisson成分を構成する。
座標$(x,y,z,a,b,c,u)$に対応するベクトル基底を$(e_0,\ldots,e_6)$とする。
\begin{equation}\label{ne:F}
\begin{aligned}
 f_{\mathrm N}&=xyz+y^3+z^3,& f_{\mathrm C}&=xz^2+y^3,\\
 F_\epsilon&=f_\epsilon+a^3+b^3+c^3+u(x^2+a^2+b^2+c^2)+u^3,
 &X_\epsilon&=V(F_\epsilon)\subset\PP^6
 \quad(\epsilon\in\{\mathrm N,\mathrm C\}).
\end{aligned}
\end{equation}
$P_1=e_0\wedge e_1\wedge e_2$、$P_2=e_3\wedge e_4\wedge e_5$、
$L=\CC P_1\oplus\CC P_2$とおく。

\begin{theorem}[非楕円型を一般点に持つ既約成分]\label{ne:main}
$X_{\mathrm N}$、$X_{\mathrm C}$は5次元の滑らかな3次超曲面である。
$P\mapsto\pi_P=(\iota_{dF_\epsilon}P)|_{X_\epsilon}$の下で
$L\simeq\mathbb A^2$は$\operatorname{Pois}(X_\epsilon)$の1つの既約成分の被約構造であり、
$P_0=P_1+P_2$でPoissonスキームそのものは滑らかである。
各成分の一般Poisson階数は4である。

$\lambda\mu\ne0$の点$\lambda P_1+\mu P_2$の3変数ブロック分解は順序を除き内在的である。
一方のブロックの平面3次曲線は、$X_{\mathrm N}$では節点型、$X_{\mathrm C}$では尖点型であり、
他方は滑らかなFermat3次曲線である。
特異なブロックを座標変換で滑らかな楕円ブロックへ取り替えることはできない。
\end{theorem}

\subsection{超曲面の滑らかさ}
\begin{lemma}\label{ne:smooth}
\eqref{ne:F}の全偏微分の共通零点は原点だけである。
\end{lemma}
\begin{proof}
$u=0$では$F_a=F_b=F_c=0$より$a=b=c=0$。
節点型の勾配$(yz,xz+3y^2,xy+3z^2)$が零なら$y=z=0$であり、
尖点型の勾配$(z^2,3y^2,2xz)$が零でも同様である。
残る$F_u=x^2$から$x=0$となる。

$u\ne0$では$u=1$に規格化する。
$a,b,c$はそれぞれ$0$または$-2/3$であり、その非零個数を$r$とすると
$a^2+b^2+c^2=4r/9$（$0\le r\le3$）。
節点型の残る3式は
\begin{equation}
 yz+2x=0,\qquad xz+3y^2=0,\qquad xy+3z^2=0.
\end{equation}
$yz=0$なら$x=y=z=0$。
$yz\ne0$なら後二式を掛けて$x^2=9yz=-18x$、従って$x=-18$。
よって$F_u=0$は、$3+4r/9=0$または$327+4r/9=0$を要求し、いずれも不可能である。
尖点型では$z^2+2x=0$、$3y^2=0$、$2xz=0$から$x=y=z=0$で、
$F_u=3+4r/9\ne0$となる。
\end{proof}

\subsection{Poissonスキームの接空間}
命題\ref{prop:pg-represent}により全双ベクトル空間は$\Lambda^3\CC^7$、次元35である。
命題\ref{prop:pg-jacobi-extension}のJacobi持上げと、勾配の1次syzygyの消滅は、
固定した$F_\epsilon$に対する線形写像と完全列を用いて示された。
これらは双対数環へのテンソル後も成立するため、Jacobi条件の線形化は
$\operatorname{Pois}(X_\epsilon)$の接空間を与える。
$Q=\sum_{i<j<k}q_{ijk}e_i\wedge e_j\wedge e_k$とする。
$Q^{ijk}$と$P_0^{ijk}$は、それぞれ$Q$と$P_0$の係数を添字の交代性で拡張したものとし、
$(F_\epsilon)_{br\ell}=\partial_b\partial_r\partial_\ell F_\epsilon$とおく。
Jacobi条件を$P_0$で線形化した整数行列$M_\epsilon$は
\begin{equation}\label{ne:M}
\begin{split}
 (M_\epsilon Q)_{ijk\alpha\ell}=\sum_{r,b}\bigl(&Q^{ir\alpha}P_0^{jkb}+P_0^{ir\alpha}Q^{jkb}
  +Q^{jr\alpha}P_0^{kib}+P_0^{jr\alpha}Q^{kib}\\
 &+Q^{kr\alpha}P_0^{ijb}+P_0^{kr\alpha}Q^{ijb}\bigr)(F_\epsilon)_{br\ell}
\end{split}
\end{equation}
で与えられる。$i<j<k$、$\alpha,\ell\in\{0,\ldots,6\}$であり、
行$(i,j,k,\alpha,\ell)$と列$(i,j,k)$を辞書式順で並べる。
零行を除くと節点型408行、尖点型312行となる。

\begin{lemma}[Jacobi方程式の線形化の階数]\label{ne:rank}
両行列の階数は33、核は$L$である。
列$(012)$と$(345)$を除いた行列には、次の行列式を持つ33次小行列
$M_{\mathrm N}'$、$M_{\mathrm C}'$が存在する：
\begin{equation}
 \det M_{\mathrm N}'=2^{24}3^{12}=8916100448256,\qquad
 \det M_{\mathrm C}'=-2^{33}3^{12}=-4565043429507072.
\end{equation}
\end{lemma}
\begin{proof}
$\lambda P_1+\mu P_2$は、中心変数$u$を共有する2つの3変数Jacobian括弧を与える。
ブロック間の括弧は零で、各ブロックの係数は他方の変数に依存しない。
よってJacobi恒等式が成り立ち、$(012)$、$(345)$の2列は零である。
従って階数は高々33であり、$L$は核に含まれる。

下界は整数上の小行列式によって得られる。
式\eqref{ne:M}から生成した$M_\epsilon$の転置を有理数上で行簡約し、
そのピボット列に対応する$M_\epsilon$の33行を選ぶ。
列$(012)$と$(345)$を除いた33次小行列の行列式を整数上で計算すると、
表示した2つの値となる。いずれも非零なので、階数は33、核は$L$である。
\end{proof}

選択した行・列の添字と小行列の全成分は、
\href{https://github.com/uedakazushi/poisson-cubic/tree/1f8d4c305a9c5c4eff4c21657254255a1a6ebf51/computations}{補足計算資料}の
\path{computations/nonelliptic_components_certificate.json}に収録している。
検算コード\path{computations/verify_nonelliptic_components.py}は
式\eqref{ne:M}から行列を再生成し、このデータと行列式の値を照合する。
この整数計算はPythonとSymPyを用い、配布ソースのルートで次を実行すれば再現できる。
\begin{quote}
\small\verb|python3 computations/verify_nonelliptic_components.py|
\end{quote}

\begin{proof}[定理\ref{ne:main}の成分性と階数]
$L$は2次元の閉既約部分スキームであり、その全点がPoissonである。
$P_0$での接空間は2次元なので、Poissonスキームの局所環はKrull次元2、
埋込み次元2の正則局所環となる。
これは$\CC$上有限型なので、$P_0$での滑らかさを意味する。
同じ次元を持つ閉部分$L$はその唯一の局所既約成分であり、局所的にスキームとして一致する。
従って$L$は大域的な既約成分で、一般点で被約である。

点$[0:1:0:-1:0:0:0]$は両$X_\epsilon$に属する。
$\lambda\mu\ne0$のとき、そこでアフィン双ベクトルの像は
$\langle e_0,e_2,e_4,e_5\rangle$で階数4。
Eulerベクトル$e_1-e_3$はその像に含まれないため、射影化しても階数4である。
ブロックの内在性には補題\ref{ne:blocks}を用いる。
\end{proof}

\begin{proof}[定理\ref{ne:main}のブロックに関する主張]
$\lambda\mu\ne0$では支持は$\langle e_0,e_1,e_2\rangle\oplus\langle e_3,e_4,e_5\rangle$。
補題\ref{ne:blocks}により2ブロックは一意である。
$F_\epsilon$の対応する射影平面への制限は$f_\epsilon$および$a^3+b^3+c^3$である。
$f_{\mathrm N}$は$[1:0:0]$近傍で$yz+y^3+z^3$となり節点、
$f_{\mathrm C}$は$z^2+y^3$となり尖点を持つ。
各式は$x$について1次で、その係数と定数項が互いに素なので既約である。
この特異性は内在的ブロックへの制限であり、座標の選択では除去できない。
\end{proof}

\subsection{特異ブロックと相対成分}
上の平面3次曲線の一般化Jacobianは次のようになる。
節点3次曲線では$\Pic^0\simeq\mathbb G_m$、尖点型では$\Pic^0\simeq\mathbb G_a$。
前者は正規化$\PP^1$上の2点のファイバーを非零スカラーで接着することで、
後者は$\CC[[t^2,t^3]]\subset\CC[[t]]$の単元商の1次係数$t$により得られる。

定理\ref{thm:pg-fermat-relative-count}により、$\lambda\mu\ne0$の点は
一般階数4の唯一の相対成分上にある。
その階数層では、ブロック制限の滑らかさはFermat式を含む非空開条件である。
従って本節の固定底の成分は、相対成分内の特異ブロックをもつ部分に位置する。
これらのブロックの量子化と行列式曲線は\cite{CubicQuantization}で扱う。
